% !TEX root = ../main.tex
% SECTION 9: WELLENOPTIK
% ============================================================

\StartChapter{Wellenoptik}

\begin{runintext}
Beschreibung von Licht als \ZSFkeyword*{elektromagnetische Welle}. Umfasst \ZSFkeyword*{Brechung}, \ZSFkeyword*{Reflexion}, \ZSFkeyword*{Polarisation} sowie wellenspezifische Phänomene wie \ZSFkeyword*{Interferenz} und \ZSFkeyword*{Gitterspektren}.
Die zugrunde liegende Wellenausbreitung ist in \ZSFref{sec:wave-equation} zusammengefasst.
\end{runintext}

\begin{tablebox}[Kenngrössen von Wellen]
\begin{ZSFtableFlat}{Y{0.4} Y{1.1} Y{1.5}}
$T$ & \textbf{Periodendauer} & Zeit bis Welle sich wiederholt \\
$\lambda$ & \textbf{Wellenlänge} & $\lambda = v \cdot T = \frac{v}{f}$ \newline Abstand zweier Wellenberge \\
$f, \nu$ & \textbf{Frequenz} & $f = \frac{1}{T} = \frac{\omega}{2\pi} = \frac{v}{\lambda}$ \newline Schwingungen pro Sekunde \\
$\omega$ & \textbf{Kreisfrequenz} & $\omega = 2\pi f = \frac{2\pi}{T}$ \newline Überstrichener Winkel pro Sekunde \\
$v, c$ & \textbf{Phasengeschw.} & $v = f \lambda = \frac{\omega}{k}$ \newline Ausbreitungsgeschw. der Welle \\
$k$ & \textbf{Wellenzahl} & $k = \frac{2\pi}{\lambda}$ \\
$\varphi$ & \textbf{Phasenversch.} & $\Delta\varphi = k\,\Delta x = \frac{2\pi}{\lambda}\Delta x$ \newline Phasenwinkel in $\si{rad}$ bei Versatz $\Delta x$ bzw. $\Delta t$ \\
\end{ZSFtableFlat}
\end{tablebox}

\SubsectionBar[sec:light-refraction-polarization]{Licht, Brechzahl, Polarisation}
\ZSFindex{Intensität}
\ZSFindex{Brechung / Brechungsgesetz}
\ZSFindex{Brechzahl / Brechungsindex}
\ZSFindex{Reflexion}
\ZSFindex{Polarisation}
\ZSFindex{Totalreflexion}
\ZSFindex{Malus'sches Gesetz}
\ZSFindexsee{Gesetz von Malus}{Malus'sches Gesetz}
\ZSFindex{Grenzwinkel}
\ZSFindex{Snellius'sches Gesetz}
\ZSFindex{Brewster-Winkel}
\begin{runintext}
Grundlagen der Lichtausbreitung in verschiedenen Medien und Verhalten des elektrischen Feldvektors (\ZSFkeyword{Polarisation}).
\end{runintext}
\begin{formulabox}[Licht im Medium: Geschwindigkeit und Wellenlänge]
\[ c = 299\,792\,458\,\si{m/s} \]
\[ c_n = \frac{c}{n} \quad \lambda_n = \frac{\lambda_0}{n} \quad \text{Brechzahl} \]
\formulanote{$c_n$ Lichtgeschwindigkeit und $\lambda_n$ Wellenlänge im Medium; $n=c/c_n$.}
\end{formulabox}
\begin{runintext}
\ZSFkeyword{Dispersion}: $n(\lambda)$; kurze Wellenlängen (blau/violett) werden stärker gebrochen als lange (rot), z.\,B.\ im Prisma.
\end{runintext}
\begin{formulabox}[Reflexion und Brechung]
\[ \theta_1' = \theta_1 \quad \text{Reflexion} \]
\[ n_1\sin\theta_1 = n_2\sin\theta_2 \quad \text{Snellius} \]
\[ I = \left(\frac{n_1-n_2}{n_1+n_2}\right)^2 I_0 \quad \text{reflekt. Intensität, senkrecht} \]
\formulanote{$n_1, n_2$ Brechzahlen (Medium 1/2); $\theta_1$ Einfallswinkel; $\theta_1'$ Ausfallswinkel; $I_0$ einfallende Intensität. Bildentstehung an gekrümmten Grenzflächen \ZSFref{sec:reflection-refraction}.}
\end{formulabox}
\begin{formulabox}[Totalreflexion und Brewster-Winkel]
\[ \sin\theta_{\mathrm{k}} = \frac{n_2}{n_1} \quad \text{Totalreflexion, $n_1>n_2$} \]
\[ \tan\theta_{\mathrm{p}} = \frac{n_2}{n_1} \quad \text{Brewster-Winkel} \]
\formulanote{Bei $\theta_{\mathrm{p}}$ stehen reflektierter und gebrochener Strahl senkrecht; der reflektierte Strahl ist vollständig linear s-polarisiert.}
\end{formulabox}
\begin{formulabox}[Polarisation: Intensität]
\[ I_1 = \frac{1}{2}I_0 \quad \text{unpolarisiert} \]
\[ I_2 = I_1\,\cos^2\alpha \quad \text{Malus} \]
\formulanote{Malus: $\alpha$ Winkel zwischen den Polarisatoren. Ein idealer erster Polarisator halbiert die Intensität unpolarisierten Lichts. Die Bedeutung der Intensität als Energiestrom steht in \ZSFref{sec:poynting-vector}.}
\end{formulabox}
\begin{statementbox}[Lineare und zirkulare Polarisation]
\ZSFkeyword{Polarisation} bezeichnet die Ausrichtung des elektrischen Feldvektors. Linear: Schwingung entlang einer festen Achse. Zirkular: Der elektrische Feldvektor rotiert um die Ausbreitungsrichtung; der Energiefluss $\vec{S}$ bleibt geradlinig. Die räumliche Orientierung von $\vec{E}$, $\vec{B}$ und Ausbreitungsrichtung ist in \ZSFref{sec:wave-equation} dargestellt.
\end{statementbox}
\begin{statementbox}[Erzeugung von Polarisation]
Durch Absorption, Streuung, Reflexion oder \ZSFkeyword{Doppelbrechung}. Bei Absorption durch einen idealen Polarisator verliert unpolarisiertes Licht die Hälfte seiner Intensität, da die Feldorientierung zufällig variiert.
\end{statementbox}

\ZSFfig[5.2cm]{graphics/Wellenoptik_Brewster_Polarisation.png}{Brewster-Winkel \& Polarisation}{Bei $\tan\theta_{\mathrm{p}} = n_2/n_1$ stehen reflektierter und gebrochener Strahl senkrecht ($90^\circ$). Reflektierter Strahl ist vollständig linear (s-)polarisiert.}

\ZSFfig{graphics/Wellenoptik_Brechung_Totalreflexion.png}{Lichtbrechung \& Totalreflexion}{Brechung bei $\theta_1 < \theta_{\mathrm{k}}$, Grenzwinkel bei $\theta_{\mathrm{k}}$ ($\theta_2 = 90^\circ$), Totalreflexion bei $\theta_1 > \theta_{\mathrm{k}}$ (Übergang optisch dicht $n_1 \to$ dünn $n_2$).}



\SubsectionBar[sec:interference-grating]{Einzelspalt, Doppelspalt, Gitter}
\ZSFindex{Intensität}
\ZSFindex{Interferenz}
\ZSFindex{Gitter / Gitterspektrum}
\ZSFindex{Beugung}
\ZSFindex{Einzelspalt}
\begin{runintext}
Am \ZSFkeyword{Einzelspalt} entsteht eine breite Beugungshülle mit zentralem Maximum. Der \ZSFkeyword{Doppelspalt} erzeugt Interferenzstreifen innerhalb dieser Hülle. Viele Spalte bilden ein \ZSFkeyword*{optisches Gitter} mit scharfen Maxima zur Trennung von Wellenlängen. Nach \ZSFkeyword[Huygens'sches Prinzip]{Huygens} ist jeder Punkt der Wellenfront Ausgangspunkt einer Elementarwelle.
\end{runintext}
\begin{formulabox}[Einzelspalt: Beugungshülle]
\[ \sin\theta_1 = \frac{\lambda}{a} \quad \text{erstes Minimum} \]
\[ a\,\sin\theta_m = m\,\lambda \quad m = 1,2,3,\ldots \quad \text{weitere Minima} \]
\[ y_1 \approx \frac{\lambda\,l}{a} \quad \text{erstes Minimum auf Schirm} \]
\[ 2y_1 \approx \frac{2\lambda\,l}{a} \quad \text{Breite zentrales Maximum} \]
\formulanote{$a$ Spaltbreite; $l$ Abstand Spalt--Schirm. Schirm weit weg oder Linse dazwischen \ZSFref{sec:lenses}: $\sin\theta\approx y/l$.}
\end{formulabox}
\begin{formulabox}[Doppelspalt: Interferenzstreifen]
\begin{longformula}
\delta = \frac{\Delta r}{\lambda}\,2\pi
= \frac{\Delta r}{\lambda}\,360^\circ
\end{longformula}
\[ \Delta r = d\,\sin\theta_m = m\,\lambda \quad \text{Max. } (m=0,1,\dots) \]
\[ \Delta r = d\,\sin\theta_m = \left(m-\frac{1}{2}\right)\lambda \quad \text{Min. } (m=1,2,\dots) \]
\[ E_0 = 2\,A_0\,\cos(\delta/2) \quad \text{Gesamtamplitude} \]
\[ I = 4\,I_0\,\cos^2(\delta/2) \quad \text{Intensität} \]
\[ y_m \approx \frac{m\,\lambda\,l}{d} \quad \text{$m$-tes Maximum auf Schirm} \]
\[ \Delta y \approx \frac{\lambda\,l}{d} \quad \text{Abstand benachbarter Maxima} \]
\formulanote{$\Delta r$ Gangunterschied; $d$ Spaltabstand; $m$ Ordnung; $\theta_m$ Beugungswinkel; $I_0$ Intensität eines Einzelspalts ($I_{\mathrm{max}}=4I_0$, $I_{\mathrm{min}}=0$). Konstruktiv: $\delta=0,2\pi,\ldots$; destruktiv: $\delta=\pi,3\pi,\ldots$.}
\end{formulabox}
\begin{formulabox}[Optisches Gitter: Spektralmaxima]
\[ \Delta r = g\,\sin\theta_m = m\,\lambda \quad \text{$m$-tes Maximum} \]
\[ \mathcal{A} = \frac{|\lambda|}{\Delta\lambda} = m\,N \quad \text{Auflösungsvermögen} \]
\formulanote{$g$ Gitterkonstante ($\neq$ Gegenstandsweite \ZSFref{sec:reflection-refraction}); $N$ beleuchtete Gitterstriche; $\mathcal{A}$ Auflösungsvermögen.}
\end{formulabox}

\ZSFmanualColumnBreak
\SubsectionBar[sec:thin-films]{Dünne Schichten}
\ZSFindex{Dünne Schicht}
\begin{runintext}
Bei Reflexion an Vorder- und Rückseite der Schicht entsteht ein \ZSFkeyword{Gangunterschied} $\Delta s=2d$. Dazu kommt pro Reflexion ein \ZSFkeyword{Phasensprung}: $\pi$ nur beim Übergang zum optisch dichteren Medium ($n_1<n_2$), sonst $0$. Durchgang ohne Sprung (\ZSFref{sec:wave-equation}).
\end{runintext}
\begin{formulabox}[Dünne Schicht: Reflexionsinterferenz]
\[ \Delta s = 2d, \quad \Delta\varphi_{\text{Weg}} = \frac{2\pi \Delta s}{\lambda_n} = \frac{4\pi nd}{\lambda_0} \]
\formulasep
\[ 2nd = m\,\lambda \quad \text{(gleiche Phasensprünge)} \]
\[ 2nd = \left(m+\tfrac{1}{2}\right)\lambda \quad \text{(ein Phasensprung $\pi$)} \]
\formulanote{$d$ Schichtdicke; $n$ und Wellenlänge im Medium \ZSFref{sec:light-refraction-polarization}. Für destruktive Interferenz Bedingungen vertauschen.}
\end{formulabox}

\ZSFfig{graphics/Wellenlehre_Reflexion_Festes_Loses_Ende.png}{Reflexion: festes/loses Ende}{Festes Ende $\equiv$ optisch dichter: $\pi$; loses Ende $\equiv$ optisch dünner: keiner}

\SubsectionBar[sec:rayleigh-criterion]{Rayleigh}
\ZSFindex{Rayleigh-Kriterium}
\ZSFindexsee{Kriterium von Rayleigh}{Rayleigh-Kriterium}
\begin{runintext}
Die Beugung an runden Öffnungen (Teleskope, Auge) limitiert das maximale optische \ZSFkeyword{Auflösungsvermögen}.
\end{runintext}
\begin{runintext}
Zwei Quellen gerade noch aufgelöst, wenn Maximum der einen im ersten Minimum der anderen.
\end{runintext}
\begin{formulabox}
\[ \alpha_{\mathrm{k}} = 1{,}22\,\frac{\lambda}{D} \quad \text{kreisförmige Öffnung $D$} \]
\formulanote{$D$ Öffnungsdurchmesser ($\neq$ Brechkraft $D=1/f$ \ZSFref{sec:lenses}); $\alpha_{\mathrm{k}}$ minimale Winkelauflösung (Rayleigh-Kriterium).}
\end{formulabox}

% ============================================================
